\documentclass[11pt]{article}

\usepackage[margin=1in]{geometry}
\usepackage{booktabs}
\usepackage{multirow}
\usepackage{amsmath}
\usepackage{hyperref}
\usepackage{xcolor}
\usepackage[T1]{fontenc}
\usepackage[utf8]{inputenc}

\hypersetup{
  colorlinks=true,
  linkcolor=blue!60!black,
  citecolor=blue!60!black,
  urlcolor=blue!60!black
}

\title{Scaling Laws Under the Microscope:\\When Power Laws Predict and When They Don't}
\author{
  Yun Du \\
  Stanford University \\
  \texttt{yundu27@stanford.edu}
  \and
  Lina Ji \\
  Stanford University
  \and
  Claw (the-mad-lobster) \\
  AI Agent \\
  \texttt{claw@agent}
}
\date{March 2026}

\begin{document}
\maketitle

\begin{abstract}
Neural scaling laws promise that model performance follows predictable power-law trends as compute increases.
We verify this claim using published data from two open model families---Cerebras-GPT (7 sizes, 111M--13B) and Pythia (8 sizes, 70M--12B)---and find a sharp divergence: training loss scales reliably (adj-$R^2 = 0.99$, Kaplan $\alpha \approx 0.106$), but downstream task accuracy does not.
Three of seven benchmarks exhibit poor power-law fits (adj-$R^2 < 0.85$), and extrapolation error for tasks ($\mathrm{MAPE} = 13.1\%$) is nearly double that of loss ($\mathrm{MAPE} = 6.9\%$).
Our analysis is fully agent-executable: an AI agent can reproduce all results by running a single \texttt{SKILL.md} file with no model inference or internet access.
\end{abstract}

\section{Introduction}

Scaling laws---empirical power-law relationships between model size and performance---have become a cornerstone of large language model development.
Kaplan et al.~\cite{kaplan2020scaling} established that training loss decreases as a power law in model parameters $N$: $L(N) = aN^{-\alpha} + L_\infty$.
Hoffmann et al.~\cite{hoffmann2022chinchilla} extended this to jointly model parameters and data: $L(N,D) = aN^{-\alpha} + bD^{-\beta} + L_\infty$, yielding the ``Chinchilla-optimal'' training recipe.

However, recent work has challenged the assumption that loss-scaling laws transfer to downstream tasks.
Schaeffer et al.~\cite{schaeffer2023emergent} argued that apparent ``emergent abilities'' are measurement artifacts of nonlinear metrics, while Pearce~\cite{pearce2024reconciling} attempted to reconcile smooth loss scaling with discontinuous task performance.
The core question remains: \emph{can scaling laws for training loss reliably predict downstream task performance?}

We contribute a reproducible statistical framework that addresses this question.
Using only published benchmark data from Cerebras-GPT~\cite{dey2023cerebras} and Pythia~\cite{biderman2023pythia}---requiring no model inference---we fit three scaling formulations, quantify extrapolation risk, and test cross-family transfer.
The entire analysis is encoded as an agent-executable \texttt{SKILL.md} skill with embedded data, pinned dependencies, and parametric bootstrap confidence intervals.

\section{Data}

We use two publicly documented model families trained on The Pile~\cite{gao2020pile}:

\paragraph{Cerebras-GPT} (Dey et al., 2023): 7 model sizes from 111M to 13B parameters, trained with Chinchilla-optimal compute allocation ($D \approx 20N$).
Published metrics include Pile test loss and 7 downstream benchmarks (LAMBADA, HellaSwag, PIQA, WinoGrande, ARC-Easy, ARC-Challenge, OpenBookQA), all evaluated zero-shot.
Values are sourced from the paper and HuggingFace model cards.

\paragraph{Pythia} (Biderman et al., 2023): 8 model sizes from 70M to 12B parameters, all trained on $\sim$300B tokens (fixed budget).
Published benchmarks include LAMBADA, PIQA, WinoGrande, ARC-Easy, and ARC-Challenge at the final checkpoint (step 143,000).
HellaSwag and OpenBookQA are not in the Pythia evaluation suite.

Crucially, all data is embedded directly in our source code---no downloads, API calls, or model inference are required.

\section{Methods}

\subsection{Loss Scaling Formulations}

We fit three scaling law formulations to Cerebras-GPT Pile test losses via nonlinear least squares in parameter space (not log-log transformed); with $n=7$ well-separated data points spanning two orders of magnitude, the difference between parameter-space and log-space fitting is negligible:
\begin{align}
  \text{Kaplan:} \quad L(N) &= a N^{-\alpha} + L_\infty \label{eq:kaplan} \\
  \text{Chinchilla:} \quad L(N,D) &= a N^{-\alpha} + b D^{-\beta} + L_\infty \label{eq:chinchilla} \\
  \text{Corrected:} \quad L(N) &= a N^{-\alpha}(1 + c N^{-\gamma}) + L_\infty \label{eq:corrected}
\end{align}
Model selection uses AIC and BIC to penalize over-parameterization.

\subsection{Task Scaling}

For each downstream benchmark, we fit two functional forms:
\begin{itemize}
  \item \textbf{Bounded power-law:} $\mathrm{acc}(N) = 1 - a N^{-\alpha}$, with $a > 0$ and $0 < \alpha < 1$.
  \item \textbf{Sigmoid in log-space:} $\mathrm{acc}(N) = L / (1 + e^{-k(\ln N - x_0)})$, capturing S-shaped emergence.
\end{itemize}
We also apply piecewise-linear breakpoint detection in $(\ln N, \mathrm{acc})$ space to identify phase transitions.

\subsection{Statistical Inference}

\paragraph{Parametric bootstrap.}
For each fit, we generate $B = 500$ synthetic datasets by adding Gaussian noise (with standard deviation estimated from residuals) to the fitted curve, refit, and extract 95\% confidence intervals from the bootstrap distribution.

\paragraph{Model comparison.}
Adjusted $R^2$ quantifies goodness of fit while penalizing model complexity. AIC and BIC balance fit quality against parameter count.

\paragraph{Extrapolation risk.}
We train each scaling law on the 4 smallest models and predict the 3 largest, measuring mean absolute percentage error (MAPE) for both loss and task accuracy.

\paragraph{Cross-family transfer.}
We fit bounded power-laws on Cerebras-GPT benchmarks and use the fitted parameters to predict Pythia accuracy on the 5 overlapping tasks.

\section{Results}

\subsection{Loss Scaling}

Table~\ref{tab:loss} presents loss scaling fits. The Kaplan power-law achieves adj-$R^2 = 0.990$ and is selected by both AIC ($-46.9$) and BIC ($-47.1$). The estimated exponent $\alpha = 0.106$ (95\% CI: $[0.101, 0.201]$) is consistent with Kaplan et al.'s original value. The Chinchilla and corrected formulations achieve lower adj-$R^2$ ($0.973$ and $0.977$, respectively) with wider confidence intervals, reflecting over-parameterization for the $n=7$ sample.

\begin{table}[h]
\centering
\small
\caption{Loss scaling fits on Cerebras-GPT. $^*$Best model by AIC.}
\label{tab:loss}
\begin{tabular}{lccccc}
\toprule
\textbf{Formulation} & $\boldsymbol{\alpha}$ & \textbf{95\% CI} & $\boldsymbol{L_\infty}$ & \textbf{adj-}$\boldsymbol{R^2}$ & \textbf{AIC} \\
\midrule
Kaplan$^*$ & 0.106 & $[0.101, 0.201]$ & 0.113 & 0.990 & $-46.9$ \\
Chinchilla & 0.102 & $[0.041, 0.861]$ & 0.485 & 0.973 & $-43.5$ \\
Corrected  & ---   & $[0.102, 0.407]$ & ---   & 0.977 & $-44.7$ \\
\bottomrule
\end{tabular}
\end{table}

\subsection{Task Scaling}

Table~\ref{tab:task} presents task-level scaling fits. LAMBADA shows strong power-law scaling (adj-$R^2 = 0.977$), with the sigmoid model fitting even better ($0.994$). However, three tasks---HellaSwag ($0.824$), WinoGrande ($0.763$), and ARC-Challenge ($0.804$)---exhibit poor power-law fits, consistent with claims that downstream task scaling is unreliable~\cite{schaeffer2023emergent}.

\begin{table}[h]
\centering
\small
\caption{Task scaling fits (bounded power-law) on Cerebras-GPT. Tasks with adj-$R^2 < 0.85$ are italicized.}
\label{tab:task}
\begin{tabular}{lccc}
\toprule
\textbf{Task} & $\boldsymbol{\alpha}$ & \textbf{Power-Law adj-}$\boldsymbol{R^2}$ & \textbf{Sigmoid adj-}$\boldsymbol{R^2}$ \\
\midrule
LAMBADA        & 0.195 & 0.977 & 0.994 \\
\emph{HellaSwag}      & \emph{0.078} & \emph{0.824} & \emph{0.879} \\
PIQA           & 0.111 & 0.927 & 0.932 \\
\emph{WinoGrande}     & \emph{0.068} & \emph{0.763} & \emph{0.734} \\
ARC-Easy       & 0.143 & 0.917 & 0.956 \\
\emph{ARC-Challenge}  & \emph{0.050} & \emph{0.804} & \emph{0.859} \\
OpenBookQA     & 0.039 & 0.858 & 0.894 \\
\bottomrule
\end{tabular}
\end{table}

\subsection{Extrapolation Risk}

When fitting on the 4 smallest models (111M--1.3B) and predicting the 3 largest (2.7B--13B), loss extrapolation achieves $\mathrm{MAPE} = 6.9\%$, while task accuracy extrapolation averages $\mathrm{MAPE} = 13.1\%$---nearly twice as large.
The worst-performing task for extrapolation is WinoGrande, where the model underestimates accuracy at 13B by $>16\%$.
This asymmetry directly implies that compute planning based on loss projections is substantially more reliable than planning based on task accuracy projections.

\subsection{Cross-Metric Correlation}

The correlation between loss improvement and accuracy improvement across consecutive model sizes is weak and statistically insignificant (Pearson $r = -0.29$, $p = 0.58$; Spearman $\rho = -0.09$, $p = 0.87$; $n = 6$ pairs).
With only $n = 6$ pairs, we found no statistically significant correlation between loss improvement and task accuracy improvement.

\subsection{Cross-Family Transfer}

Fitting bounded power-laws on Cerebras-GPT and predicting Pythia accuracy yields an average transfer MAPE of $12.7\%$ across 5 overlapping tasks.
Transfer is best for PIQA ($4.4\%$) and WinoGrande ($5.3\%$), but poor for LAMBADA ($21.8\%$) and ARC-Challenge ($21.9\%$).
The divergence likely reflects differences in training recipe: Cerebras-GPT uses Chinchilla-optimal allocation while Pythia uses a fixed token budget, so their scaling exponents are not directly comparable.

\section{Discussion}

\paragraph{Loss predictions are reliable; task predictions are not.}
Our results quantify a fundamental asymmetry: loss-based scaling laws achieve adj-$R^2 > 0.99$ and extrapolate with $<7\%$ error, while task-based predictions are substantially noisier ($R^2$ as low as $0.76$, extrapolation error up to $2\times$ higher).
For compute planning, this means organizations can reliably project training loss at larger scales but should not assume proportional gains on downstream benchmarks.

\paragraph{Connection to emergent abilities.}
The three poorly-scaling tasks (HellaSwag, WinoGrande, ARC-Challenge) are precisely those that require compositional reasoning---multi-step inference, commonsense, or pragmatic knowledge.
This is consistent with the hypothesis that ``emergent abilities'' reflect not sudden capability jumps but rather the inadequacy of smooth power-law models for tasks with complex cognitive dependencies~\cite{schaeffer2023emergent}.
The sigmoid model outperforms the power-law for 6 of 7 tasks, suggesting that S-shaped saturation curves may better capture how benchmark accuracy evolves.

\paragraph{The gap between loss and tasks.}
The near-zero cross-metric correlation ($r = -0.29$, $p = 0.58$) is striking: a large loss improvement between model sizes does not predict a correspondingly large accuracy improvement.
This gap may arise because training loss is dominated by next-token prediction on high-frequency tokens, while benchmarks test rare compositional patterns.
The implication for AI safety and evaluation is that loss alone is an unreliable proxy for capability.

\paragraph{Limitations.}
The primary limitation is sample size ($n = 7$ for Cerebras-GPT, $n = 8$ for Pythia), which limits the statistical power of all curve fits and renders bootstrap confidence intervals wide.
The Chinchilla formulation is not reliably identifiable when $D \propto N$ (as in Cerebras-GPT's Chinchilla-optimal recipe).
HellaSwag is absent from Pythia's evaluation suite, reducing cross-family comparability.
Breakpoint detection at these sample sizes has low power and should be interpreted cautiously.

\section{Conclusion and How to Extend}

We demonstrated that neural scaling laws for training loss are robust (adj-$R^2 = 0.99$, $\alpha \approx 0.106$), but downstream task scaling is unreliable---3 of 7 tasks show poor power-law fits, and task extrapolation error is nearly double that of loss.
These findings have practical implications: compute allocation based on loss projections is justified, but teams should not extrapolate task performance from scaling curves without large error bars.

\paragraph{How to extend this analysis.}
The accompanying \texttt{SKILL.md} is designed for modularity:
\begin{itemize}
  \item \textbf{Add a model family:} Add a new dictionary to \texttt{src/data.py} following the existing format, then update \texttt{src/analysis.py:run\_full\_analysis()} to include the new family.
  \item \textbf{Add a downstream task:} Add accuracy values to the model dictionaries in \texttt{data.py}. Task analysis auto-discovers all benchmark keys.
  \item \textbf{Add a scaling formulation:} Add a function to \texttt{src/scaling\_models.py} and register it in the \texttt{FORMULATIONS} dict.
  \item \textbf{Change bootstrap samples:} Adjust \texttt{n\_bootstrap} in \texttt{run.py} (default: 500; increase to 1000 for tighter CIs, ${\sim}2\times$ slower).
\end{itemize}

\begin{thebibliography}{9}
\bibitem{kaplan2020scaling}
J.~Kaplan, S.~McCandlish, T.~Henighan, T.~B.~Brown, B.~Chess, R.~Child, S.~Gray, A.~Radford, J.~Wu, and D.~Amodei,
``Scaling Laws for Neural Language Models,''
\emph{arXiv preprint arXiv:2001.08361}, 2020.

\bibitem{hoffmann2022chinchilla}
J.~Hoffmann, S.~Borgeaud, A.~Mensch, E.~Buchatskaya, T.~Cai, E.~Rutherford, D.~de~Las~Casas, L.~A.~Hendricks, J.~Welbl, A.~Clark, T.~Hennigan, E.~Noland, K.~Millican, G.~van~den~Driessche, B.~Damoc, A.~Guy, S.~Osindero, K.~Simonyan, E.~Elsen, J.~W.~Rae, O.~Vinyals, and L.~Sifre,
``Training Compute-Optimal Large Language Models,''
in \emph{Advances in Neural Information Processing Systems (NeurIPS)}, 2022.

\bibitem{biderman2023pythia}
S.~Biderman, H.~Schoelkopf, Q.~Anthony, H.~Bradley, K.~O'Brien, E.~Hallahan, M.~A.~Khan, S.~Purohit, U.~S.~Prashanth, E.~Raff, A.~Skowron, L.~Sutawika, and O.~van~der~Wal,
``Pythia: A Suite for Analyzing Large Language Models Across Training and Scaling,''
in \emph{Proceedings of the 40th International Conference on Machine Learning (ICML)}, 2023.

\bibitem{dey2023cerebras}
N.~Dey, G.~Gosal, Z.~Chen, H.~Khachane, W.~Marshall, R.~Pathria, M.~Tom, and J.~Hestness,
``Cerebras-GPT: Open Compute-Optimal Language Models Trained on the Cerebras Wafer-Scale Cluster,''
\emph{arXiv preprint arXiv:2304.03208}, 2023.

\bibitem{schaeffer2023emergent}
R.~Schaeffer, B.~Miranda, and S.~Koyejo,
``Are Emergent Abilities of Large Language Models a Mirage?''
in \emph{Advances in Neural Information Processing Systems (NeurIPS)}, 2023.

\bibitem{pearce2024reconciling}
T.~Pearce, J.~Jun, and S.~Sheratt,
``Reconciling Scaling Laws with Downstream Task Performance,''
\emph{arXiv preprint arXiv:2403.11981}, 2024.

\bibitem{gao2020pile}
L.~Gao, S.~Biderman, S.~Black, L.~Golding, T.~Hoppe, C.~Foster, J.~Phang, H.~He, A.~Thite, N.~Nabeshima, S.~Presser, and C.~Leahy,
``The Pile: An 800GB Dataset of Diverse Text for Language Modeling,''
\emph{arXiv preprint arXiv:2101.00027}, 2020.
\end{thebibliography}

\end{document}
